\documentclass[12pt]{article}
\usepackage[utf8]{inputenc}
\usepackage[T1]{fontenc}
\usepackage{hyperref}
\title{Wanggom Language Resources for K-12 Curriculum Development}
\author{Compiled by Claude Sonnet 5.5 (claude-sonnet-5-5) for Andr\'as Kornai}
\date{}
\begin{document}
\maketitle

\section{Introduction}
Wanggom (ISO 639-3 wng) is spoken in Boven Digoel, Papua, Indonesia, around Firiwage and Kawagit. Glottolog places it in the Awyu-Dumut branch of Greater Awyu (Nuclear Trans New Guinea). Its status is disputed: Ethnologue and Wikipedia treat it as a variety of Kombai, while Glottolog and Joshua Project treat it as its own language. Joshua Project gives a people population of 1,600 and says Scripture translation has started. I did not verify these at their sources.

\section{Primary Educational Resources}
No primers, textbooks or school programmes for Wanggom were found.

\section{Linguistic and Reference Materials}
Group overview: \url{https://en.wikipedia.org/wiki/Wanggom_people}. The related Kombai language article: \url{https://en.wikipedia.org/wiki/Kombai_language}. Glottolog classification: \url{https://glottolog.org/resource/languoid/id/wang1299.bigmap.html}. The Trans New Guinea database holds wordlist data, including the Kriens and Lebold (2010) survey source: \url{https://transnewguinea.org/language/wanggom?page=5} and \url{https://transnewguinea.org/source/kriens-and-lebold-2010?page=41}.

\section{Religious and Community Resources}
Joshua Project people-group profile: \url{https://joshuaproject.net/index.php/people_groups/19493}.

\section{Digital and Audio Resources}
None found.

\section{Curriculum Development Context}
Only wordlists and demographic descriptions seem to exist. A curriculum would need documentation work first, and a decision on how Wanggom relates to Kombai, where more materials may exist. Indonesian is the school language in the region.

\end{document}